% Shared preamble for learn-attn documents
% Usage: % Shared preamble for learn-attn documents
% Usage: % Shared preamble for learn-attn documents
% Usage: \input{preamble} at the top of each document

\documentclass[11pt, a4paper]{article}

\usepackage[T1]{fontenc}
\usepackage{lmodern}
\usepackage{microtype}
\usepackage[margin=1in]{geometry}
\usepackage{amsmath, amssymb}
\usepackage{booktabs}
\usepackage{enumitem}
\usepackage{parskip}
\usepackage[dvipsnames]{xcolor}
\usepackage{listings}
\usepackage{courier}
\usepackage{pgfplots}
\pgfplotsset{compat=1.18}
\usetikzlibrary{backgrounds}
\usepackage[colorlinks=true, linkcolor=NavyBlue, urlcolor=NavyBlue, citecolor=NavyBlue]{hyperref}

% Code listing style
\definecolor{codebg}{gray}{0.96}
\definecolor{codeframe}{gray}{0.8}
\definecolor{codecomment}{RGB}{90,130,90}
\definecolor{codestring}{RGB}{160,60,60}
\definecolor{codekw}{RGB}{40,40,120}

\lstset{
  language=Python,
  basicstyle=\ttfamily\small,
  backgroundcolor=\color{codebg},
  frame=single,
  rulecolor=\color{codeframe},
  framesep=6pt,
  xleftmargin=6pt,
  xrightmargin=6pt,
  breaklines=true,
  breakatwhitespace=true,
  showstringspaces=false,
  keywordstyle=\color{codekw}\bfseries,
  commentstyle=\color{codecomment}\itshape,
  stringstyle=\color{codestring},
  numbers=none,
  tabsize=4,
  aboveskip=1em,
  belowskip=1em,
  morekeywords={self, None, True, False},
}

% Inline code
\newcommand{\code}[1]{\texttt{#1}}

% Document series header
\newcommand{\docheader}[2]{%
  \begin{center}
    {\Large\bfseries #1}\\[6pt]
    {\itshape Document #2 of 6\,---\,Building a GPT from Scratch}
  \end{center}
  \vspace{1em}
}

% --- Code block annotations ---
% Small colored tags placed before a listing to clarify its role.

% Code that the reader should put into a project file.
% Usage: \filetag{src/model.py}  (before \begin{lstlisting})
\newcommand{\filetag}[1]{%
  \par\noindent
  \colorbox{black!8}{\small\sffamily\bfseries File: \code{#1}}%
  \nopagebreak\par\nopagebreak\vspace{-0.6em}%
}

% Standalone demo — try in a Python REPL or scratch script, not part of the project.
\newcommand{\demotag}{%
  \par\noindent
  \colorbox{NavyBlue!12}{\small\sffamily\bfseries Demo \textnormal{--- try in a REPL}}%
  \nopagebreak\par\nopagebreak\vspace{-0.6em}%
}

% Pseudocode or conceptual snippet illustrating an idea.
\newcommand{\concepttag}{%
  \par\noindent
  \colorbox{black!6}{\small\sffamily\bfseries Conceptual \textnormal{--- illustrating the idea}}%
  \nopagebreak\par\nopagebreak\vspace{-0.6em}%
}

% Expected terminal or REPL output.
\newcommand{\outputtag}{%
  \par\noindent
  \colorbox{black!6}{\small\sffamily\bfseries Expected output}%
  \nopagebreak\par\nopagebreak\vspace{-0.6em}%
}

% Shell command the reader should run.
% Usage: \shelltag  (before a listing with language=bash or similar)
\newcommand{\shelltag}{%
  \par\noindent
  \colorbox{Goldenrod!20}{\small\sffamily\bfseries Shell \textnormal{--- run in your terminal}}%
  \nopagebreak\par\nopagebreak\vspace{-0.6em}%
}

% Verification checkpoint — run this and compare.
\newcommand{\checkpoint}[1]{%
  \medskip\par\noindent
  \fbox{\parbox{\dimexpr\linewidth-2\fboxsep-2\fboxrule}{%
    {\sffamily\bfseries Checkpoint:} #1}}%
  \medskip
}

% Tip — clarifies a common point of confusion.
\newcommand{\tip}[1]{%
  \medskip\par\noindent
  \colorbox{ForestGreen!10}{\parbox{\dimexpr\linewidth-2\fboxsep}{%
    {\sffamily\bfseries\color{ForestGreen!70!black} Tip:} #1}}%
  \medskip
}
 at the top of each document

\documentclass[11pt, a4paper]{article}

\usepackage[T1]{fontenc}
\usepackage{lmodern}
\usepackage{microtype}
\usepackage[margin=1in]{geometry}
\usepackage{amsmath, amssymb}
\usepackage{booktabs}
\usepackage{enumitem}
\usepackage{parskip}
\usepackage[dvipsnames]{xcolor}
\usepackage{listings}
\usepackage{courier}
\usepackage{pgfplots}
\pgfplotsset{compat=1.18}
\usetikzlibrary{backgrounds}
\usepackage[colorlinks=true, linkcolor=NavyBlue, urlcolor=NavyBlue, citecolor=NavyBlue]{hyperref}

% Code listing style
\definecolor{codebg}{gray}{0.96}
\definecolor{codeframe}{gray}{0.8}
\definecolor{codecomment}{RGB}{90,130,90}
\definecolor{codestring}{RGB}{160,60,60}
\definecolor{codekw}{RGB}{40,40,120}

\lstset{
  language=Python,
  basicstyle=\ttfamily\small,
  backgroundcolor=\color{codebg},
  frame=single,
  rulecolor=\color{codeframe},
  framesep=6pt,
  xleftmargin=6pt,
  xrightmargin=6pt,
  breaklines=true,
  breakatwhitespace=true,
  showstringspaces=false,
  keywordstyle=\color{codekw}\bfseries,
  commentstyle=\color{codecomment}\itshape,
  stringstyle=\color{codestring},
  numbers=none,
  tabsize=4,
  aboveskip=1em,
  belowskip=1em,
  morekeywords={self, None, True, False},
}

% Inline code
\newcommand{\code}[1]{\texttt{#1}}

% Document series header
\newcommand{\docheader}[2]{%
  \begin{center}
    {\Large\bfseries #1}\\[6pt]
    {\itshape Document #2 of 6\,---\,Building a GPT from Scratch}
  \end{center}
  \vspace{1em}
}

% --- Code block annotations ---
% Small colored tags placed before a listing to clarify its role.

% Code that the reader should put into a project file.
% Usage: \filetag{src/model.py}  (before \begin{lstlisting})
\newcommand{\filetag}[1]{%
  \par\noindent
  \colorbox{black!8}{\small\sffamily\bfseries File: \code{#1}}%
  \nopagebreak\par\nopagebreak\vspace{-0.6em}%
}

% Standalone demo — try in a Python REPL or scratch script, not part of the project.
\newcommand{\demotag}{%
  \par\noindent
  \colorbox{NavyBlue!12}{\small\sffamily\bfseries Demo \textnormal{--- try in a REPL}}%
  \nopagebreak\par\nopagebreak\vspace{-0.6em}%
}

% Pseudocode or conceptual snippet illustrating an idea.
\newcommand{\concepttag}{%
  \par\noindent
  \colorbox{black!6}{\small\sffamily\bfseries Conceptual \textnormal{--- illustrating the idea}}%
  \nopagebreak\par\nopagebreak\vspace{-0.6em}%
}

% Expected terminal or REPL output.
\newcommand{\outputtag}{%
  \par\noindent
  \colorbox{black!6}{\small\sffamily\bfseries Expected output}%
  \nopagebreak\par\nopagebreak\vspace{-0.6em}%
}

% Shell command the reader should run.
% Usage: \shelltag  (before a listing with language=bash or similar)
\newcommand{\shelltag}{%
  \par\noindent
  \colorbox{Goldenrod!20}{\small\sffamily\bfseries Shell \textnormal{--- run in your terminal}}%
  \nopagebreak\par\nopagebreak\vspace{-0.6em}%
}

% Verification checkpoint — run this and compare.
\newcommand{\checkpoint}[1]{%
  \medskip\par\noindent
  \fbox{\parbox{\dimexpr\linewidth-2\fboxsep-2\fboxrule}{%
    {\sffamily\bfseries Checkpoint:} #1}}%
  \medskip
}

% Tip — clarifies a common point of confusion.
\newcommand{\tip}[1]{%
  \medskip\par\noindent
  \colorbox{ForestGreen!10}{\parbox{\dimexpr\linewidth-2\fboxsep}{%
    {\sffamily\bfseries\color{ForestGreen!70!black} Tip:} #1}}%
  \medskip
}
 at the top of each document

\documentclass[11pt, a4paper]{article}

\usepackage[T1]{fontenc}
\usepackage{lmodern}
\usepackage{microtype}
\usepackage[margin=1in]{geometry}
\usepackage{amsmath, amssymb}
\usepackage{booktabs}
\usepackage{enumitem}
\usepackage{parskip}
\usepackage[dvipsnames]{xcolor}
\usepackage{listings}
\usepackage{courier}
\usepackage{pgfplots}
\pgfplotsset{compat=1.18}
\usetikzlibrary{backgrounds}
\usepackage[colorlinks=true, linkcolor=NavyBlue, urlcolor=NavyBlue, citecolor=NavyBlue]{hyperref}

% Code listing style
\definecolor{codebg}{gray}{0.96}
\definecolor{codeframe}{gray}{0.8}
\definecolor{codecomment}{RGB}{90,130,90}
\definecolor{codestring}{RGB}{160,60,60}
\definecolor{codekw}{RGB}{40,40,120}

\lstset{
  language=Python,
  basicstyle=\ttfamily\small,
  backgroundcolor=\color{codebg},
  frame=single,
  rulecolor=\color{codeframe},
  framesep=6pt,
  xleftmargin=6pt,
  xrightmargin=6pt,
  breaklines=true,
  breakatwhitespace=true,
  showstringspaces=false,
  keywordstyle=\color{codekw}\bfseries,
  commentstyle=\color{codecomment}\itshape,
  stringstyle=\color{codestring},
  numbers=none,
  tabsize=4,
  aboveskip=1em,
  belowskip=1em,
  morekeywords={self, None, True, False},
}

% Inline code
\newcommand{\code}[1]{\texttt{#1}}

% Document series header
\newcommand{\docheader}[2]{%
  \begin{center}
    {\Large\bfseries #1}\\[6pt]
    {\itshape Document #2 of 6\,---\,Building a GPT from Scratch}
  \end{center}
  \vspace{1em}
}

% --- Code block annotations ---
% Small colored tags placed before a listing to clarify its role.

% Code that the reader should put into a project file.
% Usage: \filetag{src/model.py}  (before \begin{lstlisting})
\newcommand{\filetag}[1]{%
  \par\noindent
  \colorbox{black!8}{\small\sffamily\bfseries File: \code{#1}}%
  \nopagebreak\par\nopagebreak\vspace{-0.6em}%
}

% Standalone demo — try in a Python REPL or scratch script, not part of the project.
\newcommand{\demotag}{%
  \par\noindent
  \colorbox{NavyBlue!12}{\small\sffamily\bfseries Demo \textnormal{--- try in a REPL}}%
  \nopagebreak\par\nopagebreak\vspace{-0.6em}%
}

% Pseudocode or conceptual snippet illustrating an idea.
\newcommand{\concepttag}{%
  \par\noindent
  \colorbox{black!6}{\small\sffamily\bfseries Conceptual \textnormal{--- illustrating the idea}}%
  \nopagebreak\par\nopagebreak\vspace{-0.6em}%
}

% Expected terminal or REPL output.
\newcommand{\outputtag}{%
  \par\noindent
  \colorbox{black!6}{\small\sffamily\bfseries Expected output}%
  \nopagebreak\par\nopagebreak\vspace{-0.6em}%
}

% Shell command the reader should run.
% Usage: \shelltag  (before a listing with language=bash or similar)
\newcommand{\shelltag}{%
  \par\noindent
  \colorbox{Goldenrod!20}{\small\sffamily\bfseries Shell \textnormal{--- run in your terminal}}%
  \nopagebreak\par\nopagebreak\vspace{-0.6em}%
}

% Verification checkpoint — run this and compare.
\newcommand{\checkpoint}[1]{%
  \medskip\par\noindent
  \fbox{\parbox{\dimexpr\linewidth-2\fboxsep-2\fboxrule}{%
    {\sffamily\bfseries Checkpoint:} #1}}%
  \medskip
}

% Tip — clarifies a common point of confusion.
\newcommand{\tip}[1]{%
  \medskip\par\noindent
  \colorbox{ForestGreen!10}{\parbox{\dimexpr\linewidth-2\fboxsep}{%
    {\sffamily\bfseries\color{ForestGreen!70!black} Tip:} #1}}%
  \medskip
}

\begin{document}
\docheader{The Transformer Block}{2}

Document 01 built up the core operation of a transformer: scaled dot-product attention, multi-head attention, and causal masking. But attention alone is not enough to build a deep network. This document covers the remaining components that, together with attention, form the \textbf{transformer block} --- the repeating unit that gets stacked to build every modern large language model.

We will build four things: residual connections, layer normalization, the position-wise feed-forward network, and then assemble them into a single block.

\section{Residual Connections}

\subsection{The degradation problem}

A naive approach to building a more powerful network is to stack more layers. In theory, a deeper network can represent everything a shallower one can (the extra layers could just learn the identity function). In practice, He et al.\ (2015) showed that this does not happen: beyond a certain depth, adding layers actually \textit{increases} training error. This is not overfitting --- it is an optimization failure. The network is harder to train, and gradient-based optimization gets stuck.

The root cause is that gradients must pass through many multiplicative layers during backpropagation. Each layer multiplies the gradient by its Jacobian. If those Jacobians have eigenvalues consistently less than 1, gradients shrink exponentially (vanishing gradients). If greater than 1, they explode. Either way, the optimizer receives a poor signal about how to update early layers.

\subsection{The ResNet fix}

Instead of asking a sublayer to learn a target function $H(x)$, we ask it to learn the \textit{residual} $F(x) = H(x) - x$:

\[
\text{output} = x + F(x)
\]

This is a residual connection (also called a skip connection). The sublayer only needs to learn what to \textit{add} to the input, not how to reconstruct it from scratch.

Why does this help gradient flow? Take the gradient of the output with respect to the input:

\[
\frac{\partial}{\partial x}\bigl[x + F(x)\bigr] = I + \frac{\partial F}{\partial x}
\]

The identity matrix $I$ guarantees that even if $\partial F / \partial x$ is very small, the gradient is at least $I$. In a deep stack of $N$ blocks without residuals, the gradient must pass through $N$ multiplicative factors --- it can vanish or explode. With residuals, there is always a ``shortcut'' path where the gradient is 1. Information from the loss can reach the earliest layers without being distorted by intervening computations.

In code, a residual connection is one line:

\begin{lstlisting}
x = x + sublayer(x)
\end{lstlisting}

Every sublayer in a transformer (attention and feed-forward) is wrapped with a residual connection.

It is useful to think of the tensor flowing through the network as a \textbf{residual stream}. Each sublayer reads from the stream, computes a delta, and writes (adds) that delta back. The stream accumulates contributions from every sublayer, and the residual connections ensure that no sublayer can destroy what previous sublayers have written --- it can only add to it.

\section{Layer Normalization}

Deep networks train better when each layer's input has a stable distribution. Normalization layers address this by standardizing activations.

\subsection{Why not batch normalization?}

Batch normalization (Ioffe \& Szegedy, 2015) computes the mean and variance of each feature across the batch dimension. It works well for CNNs, but it is a poor fit for sequence models:

\begin{enumerate}
  \item \textbf{Variable-length sequences.} Different samples in a batch may have different lengths. Computing batch statistics over padded positions mixes real data with padding.
  \item \textbf{Position-dependent distributions.} Position 1 of every sequence in a batch may have a very different distribution than position 50. Averaging across the batch conflates these.
  \item \textbf{Inference with batch\_size=1.} During generation we process one sequence at a time. Batch statistics over a single example are meaningless.
\end{enumerate}

\subsection{Layer normalization}

Layer normalization (Ba et al., 2016) sidesteps all of these problems by normalizing across the \textbf{feature dimension} for each individual position, independently.

Given an input vector $\mathbf{x}$ of shape $(\ldots, d)$:

\begin{enumerate}
  \item Compute the mean across features:
\[
\mu = \frac{1}{d}\sum_{i=1}^{d} x_i
\]

  \item Compute the variance across features:
\[
\sigma^2 = \frac{1}{d}\sum_{i=1}^{d}(x_i - \mu)^2
\]

  \item Normalize, scale, and shift:
\[
\mathrm{LayerNorm}(\mathbf{x}) = \gamma \odot \frac{\mathbf{x} - \mu}{\sqrt{\sigma^2 + \varepsilon}} + \beta
\]
\end{enumerate}

where:
\begin{itemize}
  \item $\gamma$ (weight) and $\beta$ (bias) are learnable parameters of shape $(d,)$.
  \item $\varepsilon$ is a small constant (typically $10^{-5}$) to prevent division by zero.
\end{itemize}

The normalization centers and scales each position's feature vector to zero mean and unit variance. The learnable $\gamma$ and $\beta$ then let the model undo this if needed. Each position in the sequence is normalized independently, using only its own $d$-dimensional feature vector.

\subsection{Implementation}

Following modern convention (GPT-J, LLaMA), we default to \code{bias=False}:

\begin{lstlisting}
class LayerNorm(nn.Module):
    def __init__(self, n_embd, bias=False, eps=1e-5):
        super().__init__()
        self.weight = nn.Parameter(torch.ones(n_embd))
        self.bias = nn.Parameter(torch.zeros(n_embd)) if bias else None
        self.eps = eps

    def forward(self, x):
        return F.layer_norm(x, self.weight.shape, self.weight, self.bias, self.eps)
\end{lstlisting}

\section{Pre-Norm vs Post-Norm}

Where you place the LayerNorm relative to the sublayer and the residual connection matters for training stability.

\textbf{Post-norm} (the original ``Attention Is All You Need'' paper):
\[
x = \mathrm{LayerNorm}(x + \mathrm{sublayer}(x))
\]

\textbf{Pre-norm} (GPT-2, GPT-3, and most modern transformers):
\[
x = x + \mathrm{sublayer}(\mathrm{LayerNorm}(x))
\]

\subsection{Why pre-norm won}

In post-norm, the normalization happens after the residual addition. During early training, the residual $x$ and the sublayer output can have very different magnitudes. Their sum may have large variance, and the normalization must correct for this mismatch at every layer.

In pre-norm, the sublayer always receives a well-conditioned input (zero mean, unit variance) regardless of what earlier layers have done. This decouples the sublayer's optimization from the scale of the residual stream. Training is more stable, especially for deep models.

One consequence of pre-norm: the output of the last transformer block is a residual addition, but there is no normalization applied after that addition. So we need a \textbf{final LayerNorm} after the entire stack of blocks.

We use pre-norm throughout BabyGPT.

\begin{tabular}{llll}
\toprule
 & Post-norm & Pre-norm \\
\midrule
Formula & $x = \mathrm{LN}(x + \mathrm{sublayer}(x))$ & $x = x + \mathrm{sublayer}(\mathrm{LN}(x))$ \\
Used by & Original transformer & GPT-2, GPT-3, LLaMA \\
Final LN needed? & No & Yes \\
Training stability & Requires careful warmup & More stable out of the box \\
\bottomrule
\end{tabular}

\section{Position-wise Feed-Forward Network}

Attention is the inter-position operation: it lets each position gather information from other positions. The feed-forward network (FFN) is the intra-position operation: it processes each position's representation independently, with a larger capacity.

\subsection{Architecture}

The FFN is a simple two-layer MLP applied to each position:

\[
\mathrm{FFN}(\mathbf{x}) = W_2 \cdot \mathrm{GELU}(W_1 \mathbf{x} + \mathbf{b}_1) + \mathbf{b}_2
\]

where $W_1$ has shape $(d_{\text{model}}, d_{\text{ff}})$ with $d_{\text{ff}} = 4 \cdot d_{\text{model}}$, and $W_2$ has shape $(d_{\text{ff}}, d_{\text{model}})$.

\subsection{The 4x expansion}

The FFN temporarily expands to $4 \times d_{\text{model}}$ dimensions. This gives each position a much wider intermediate representation --- more capacity to extract and transform features. The expansion factor 4 is an empirical choice from the original paper that has become standard.

\subsection{Activation function}

The original transformer uses ReLU: $\max(0, x)$. GPT-2 and most modern language models use GELU:

\[
\mathrm{GELU}(x) = x \cdot \Phi(x)
\]

where $\Phi(x)$ is the standard Gaussian CDF. Unlike ReLU, which has a hard cutoff at zero, GELU is smooth everywhere --- negative inputs are dampened rather than zeroed out.

\subsection{Why ``position-wise''}

The same weights $W_1$ and $W_2$ are shared across all positions in the sequence. Each position is processed independently --- there is no information flow between positions inside the FFN. It is equivalent to a $1 \times 1$ convolution in CNN terminology.

\subsection{Implementation}

\begin{lstlisting}
class FeedForward(nn.Module):
    def __init__(self, config):
        super().__init__()
        self.c_fc = nn.Linear(config.n_embd, 4 * config.n_embd, bias=config.bias)
        self.c_proj = nn.Linear(4 * config.n_embd, config.n_embd, bias=config.bias)
        self.dropout = nn.Dropout(config.dropout)

    def forward(self, x):
        x = F.gelu(self.c_fc(x))
        return self.dropout(self.c_proj(x))
\end{lstlisting}

\section{Assembling the Transformer Block}

A single transformer block combines everything:

\begin{enumerate}
  \item Attention sublayer with residual and pre-norm: $x = x + \mathrm{CausalSelfAttention}(\mathrm{LayerNorm}(x))$
  \item FFN sublayer with residual and pre-norm: $x = x + \mathrm{FeedForward}(\mathrm{LayerNorm}(x))$
\end{enumerate}

The two sublayers serve complementary roles:
\begin{itemize}
  \item \textbf{Attention} handles inter-position communication.
  \item \textbf{FFN} handles intra-position computation.
\end{itemize}

\subsection{Implementation}

\begin{lstlisting}
class TransformerBlock(nn.Module):
    def __init__(self, config):
        super().__init__()
        self.ln_1 = LayerNorm(config.n_embd, bias=config.bias)
        self.attn = CausalSelfAttention(config)
        self.ln_2 = LayerNorm(config.n_embd, bias=config.bias)
        self.ffn = FeedForward(config)

    def forward(self, x):
        x = x + self.attn(self.ln_1(x))
        x = x + self.ffn(self.ln_2(x))
        return x
\end{lstlisting}

\section{Stacking Blocks}

A transformer model is $N$ identical blocks applied in sequence. In our BabyGPT, $N = 6$. Each block has its own parameters (not shared).

After the last block, a final LayerNorm is applied:

\begin{lstlisting}
x = tok_emb + pos_emb
for block in self.blocks:
    x = block(x)
x = self.final_ln(x)
\end{lstlisting}

\subsection{Parameter count}

With our config ($d_{\text{model}} = 384$, $d_{\text{ff}} = 1536$, $h = 6$, bias = False):

\begin{itemize}
  \item Attention: $3 \times (384 \times 384) + 384 \times 384 = 589{,}824$
  \item FFN: $384 \times 1536 + 1536 \times 384 = 1{,}179{,}648$
  \item LayerNorm: $2 \times 384 = 768$
  \item Total per block: $1{,}770{,}240$ ($\sim$1.77M parameters)
  \item 6 blocks: $\sim$10.6M parameters
\end{itemize}

The FFN accounts for about two-thirds of each block's parameters.

\section{Looking Ahead}

We now have the complete building block. In Document 03, we look at the full encoder-decoder architecture from the original paper. In Document 04, we simplify to the decoder-only GPT architecture we implement.

\end{document}
